\documentclass[../../main2.tex]{subfiles}

\begin{document}

\chapter{相关与因果：谁是真正的推手}
\label{ch:causation}

\begin{motivation}
	上一章遗留的问题：干预预测离开因果寸步难行，因果箭头只许从因推果。但现实里摆在我面前的永远是\textbf{数据}——两条曲线一起涨，凭什么说谁推了谁？「相关不等于原因」人人会背，背完照样在坑里。\\
	\textbf{核心动机}：把这句口号变成手艺——认出三种骗人的相关，拿到三把逼近因果的镊子。
\end{motivation}

\begin{logicmap}
\begin{tikzpicture}[node distance=0.85cm, every node/.style={draw, rectangle, rounded corners, align=center, font=\small}]
\node (q) {两条曲线一起涨，谁推了谁？};
\node (c) [below=of q] {第三只手：混杂};
\node (d) [below=of c] {自己骗自己：选择与确认偏差};
\node (p) [below=of d] {三把镊子：逼近因果};
\node (n) [below=of p] {分账怎么做？（抛给 Part III）};
\draw[-{Stealth}] (q) -- (c);
\draw[-{Stealth}] (c) -- (d);
\draw[-{Stealth}] (d) -- (p);
\draw[-{Stealth}] (p) -- (n);
\end{tikzpicture}
\end{logicmap}

\section{第三只手：混杂}

\begin{readingnote}
	你有没有想过，为什么「冰淇淋销量」和「溺水人数」两条曲线像双胞胎？\\
	卖冰淇淋会导致淹死人吗？还是淹死人会让人想吃冰淇淋？
\end{readingnote}

答案你已经知道了：\textbf{天热}。天热，人既买冰淇淋，也下水。两条曲线一起涨，因为它们被同一只手推着。这只手叫\textbf{混杂因子}（confounder）——一个藏在背后、同时驱动两个变量的东西。

混杂的坑，古今都在挖：

\begin{itemize}
	\item \textbf{古}：古代欧洲「鹳鸟多的村庄新生儿多」。鹳鸟送子？其实是农村——村庄既招鹳鸟，也多生娃。
	\item \textbf{今}：「喝咖啡的人更长寿」。是咖啡续命，还是喝得起咖啡的人本来就更有钱、体检更多？
	\item \textbf{今}：「上补习班的学生成绩好」。补习有用，还是本来就爱学习的孩子更会被送去补习？
	\item \textbf{今}：「网红店排队越长越好吃」。真好吃，还是花钱雇的队排出了「好吃」的\textbf{相关}？
\end{itemize}

注意最后一例：有的混杂是\textbf{人故意造的}。你看到的世界，可能有人在精心维护「相关」。

\begin{questionbox}
	\textit{现在你可能想反驳：「那把混杂因子都测出来，控制住不就行了？」}\\
	\textbf{方向对，但有个坑：你不可能控制住所有东西。}「学习时长」能测，「学习状态」难测；「家境」能测，「家里书架上有多少书」难测。漏掉的那个混杂，照样污染结论。所以控制变量是\textbf{逼近}，不是\textbf{到达}——下一节的镊子比它硬。
\end{questionbox}

\paragraph*{逻辑衔接}
	混杂是「第三只手」。但比外部混杂更麻烦的，是自己人骗自己人——下一节看两台自欺的机器。

\section{自己骗自己：偏差从哪来}

\begin{readingnote}
	你有没有想过，为什么网上待得越久，越觉得自己代表「大多数」？\\
	你看到的「大多数」，是被谁挑出来喂给你的？
\end{readingnote}

两台机器，一台外、一台内。

\textbf{外：选择偏差}。短视频平台的算法为什么专推你爱看的？别骂它坏——它的动机写在合同里：你多停一秒，它多赚一分。推荐「你爱看的」是它的最优策略，不是你的。于是你刷到的世界，是一个以「你的点击」为筛子筛过的世界：

\begin{itemize}
	\item 你觉得某明星全网都在夸——因为夸的被推给你，骂的你没刷到；
	\item 你觉得某个养生法全网都在信——因为信的人聚在一个评论区；
	\item 你觉得房价「所有人都盼着跌」——因为同温层替你把另一批人屏蔽了。
\end{itemize}

这就是信息茧房（filter bubble）的白话版：不是你蠢，是喂料的动机和你的动机不一致。平台要的是停留时长，你要的是真实世界——两个目标分岔的地方，就是偏差的产地。

\textbf{内：确认偏差}。人会主动找支持自己的证据。买了只股票，你突然刷到的全是利好；想跳槽，你突然觉得老板处处刁难。数据没变，筛子换了。

\begin{tcolorbox}[colback=green!5, title=\textit{生活类比}]
	确认偏差像只挑食的猫：碗里十条鱼，它只吃「证明我没错」那几条，然后宣称「看，全是好鱼」。混杂是别人往碗里塞了坏鱼；确认偏差是自己挑着吃。两件事叠加，你就活在一个量身定做的假世界里。
\end{tcolorbox}

流行病谣言是最好的观察窗：「板蓝根防病毒」「熏醋消毒」——每个谣言都带着一两个「真的有人好了」的案例（选择偏差），再配上「我二舅就是这么好的」（确认偏差）。案例是真的，结论是假的。

\begin{questionbox}
	\textit{现在你可能想反驳：「照你这么说，统计数据根本不可信，什么都别信了。」}\\
	\textbf{恰恰相反——知道了筛子的存在，你才能修筛子。}统计不可信的年代，人信的是「我二舅」；统计可信的前提，是你知道它可能被谁污染、怎么查。悲观不是结论，是工作的开始。
\end{questionbox}

\paragraph*{逻辑衔接}
	两台自欺机器拆开了。那么硬问题来了：一堆不可信的相关摆在桌上，怎么把因果逼出来？下一节发镊子。

\section{三把镊子：把因果逼出来}

\begin{readingnote}
	你有没有想过，「吸烟导致肺癌」当年是怎么从一堆相关里定案的？\\
	靠的不是更大的数据，是几把\textbf{能动手}的镊子。
\end{readingnote}

\textbf{镊子一：自然实验（natural experiment）}。人为实验做不了的地方，找世界「顺手做过的实验」。例：有的州涨了烟草税，有的没涨——税率像随机分组，对比两州的吸烟率变化。再如双胞胎研究：同卵双胞胎一个吸烟一个不吸，把基因这个混杂因子按住了。

\textbf{镊子二：剂量反应}。如果真的是原因，通常「越多越明显」。吸烟越多、年限越长，肺癌率越高——这条梯度让「体质论」很难看。反之，「戴不戴口罩都一样」的宣称，先问一句：剂量呢？

\textbf{镊子三：机制}。尼古丁代谢、致癌物、细胞病变——机制链一节节被找到，因果就从统计升级为物理。塞麦尔维斯当年只有干预有效、没有机制，说服力就差一口气。

三把镊子的次序，正是第 6 章「从因推果」的落地：找一个「世界替你分了组」的地方，比较两组的结果。所有镊子都在偷同一件事——\textbf{造一个平行宇宙来对比}。

\begin{tcolorbox}[colback=green!5, title=\textit{生活类比}]
	试衣服要照镜子，检验因果要照\textbf{反事实}的镜子：镜子里是没有你的那个世界。自然实验是找到一面现成的镜子；剂量反应是把镜子从远到近摆一排；机制是弄清镜子的成像原理。
\end{tcolorbox}

\begin{warningbox}
	\textbf{这套镊子什么时候失灵？}历史研究常有遗憾：没有自然实验（拿破仑只有一个）、剂量不明（古代赋税数据粗糙）、机制没找到（中医很多经验方）。此时的正确姿势不是硬下结论，而是\lowfreq{标注把握，留待后人}。一本讲预测的书，先要承认自己够不着的地方。
\end{warningbox}

\paragraph*{逻辑衔接}
	镊子能把「谁是推手」找出来。但现实里推手从来不止一个——两个原因搅在一起，「各占多少」怎么算？这个要数字的问题，科学界至今含糊，法庭却躲不掉。下一章进法庭。

\section{Part II 尾声：第二枚追问}

动机看清了，因果的手艺也拿到了。可是把它们放回社会，一个更大的裂缝裂开了：菜市场里每个摊贩都只算自己的账，菜价却有全城的脾气；千万人各自想赚钱，房价走出的却是一条谁也没画过的曲线。

\begin{warningbox}
	这是本书第二枚定时的追问：\textbf{即使搞清了每个人的动机和因果，个体加总起来，会得到什么？}\\
	还是不答。Part III 先把「分账」的手艺备齐——因为加总的第一步，就是承认每个个体\textbf{各占多少}——Part IV 开头（第 10 章）把两枚追问一起收掉。成对出现，不许只埋不收。
\end{warningbox}

\begin{reader_anticipation}
	\item 三把镊子只能「逼近」因果，那「各占多少」的精确账谁在算？→ 第 8 章进法庭：一群人被逼着算了上千年。
	\item 两枚追问一直悬着不难受吗？→ 忍到第 10 章，那里把它们一起收掉——成对出现，是全书的规矩。
\end{reader_anticipation}

\paragraph*{逻辑衔接}
	工具三件套齐了：动机（Part I）、因果（Part II）、外加两枚埋好的追问。但「各占多少」到现在还是一句嘴上话。下一章看一群人怎么被逼着把它算了上千年——欢迎来到法庭。

\begin{thinkbox}
\begin{enumerate}
	\item 翻出你最近深信的一个「相关」（比如「我家孩子玩手机所以成绩降了」）：列出至少两个可能的混杂因子，再从三把镊子里挑一把，设计一个能逼近因果的小检验。
	\item \textbf{反事实}：如果短视频平台的算法目标从「停留时长」改成「用户三个月后的满意度」，你的信息茧房会变薄还是变厚？平台的动机变了，喂给你的世界会怎么变？
	\item \textbf{反事实}：如果 1950 年代烟草公司资助的研究「证明」了体质论——吸烟和肺癌只是同根——今天的医学教科书会少哪一章？「自然实验」这把镊子，会晚多少年才被重视？
\end{enumerate}
\end{thinkbox}

\end{document}
